\documentclass[a4paper,11pt]{article}
\usepackage{exam-style-341}
\examinehead{341 农业综合知识三 · 程序设计}{模拟试卷（二）}
\begin{document}
\examcover{341 农业综合知识三}{程序设计部分 · 模拟试卷（二）}{程序设计（C 语言）（50 分）}{50 分}{60 分钟}{本科目只考程序设计，与数据库、网络无关。}

\parttitle{程序设计（C 语言）}

\qtypename{一、单项选择题}{10 题，每题 2 分，共 20 分}
\begin{qgroup}
\q 设有定义 \code{int a = 5, b = 2; float x = 3.0;}，则表达式 \code{a / b * x} 的值是（\quad）。
\choicesline{A. 6.0}{B. 7.5}{C. 2.5}{D. 15.0}

\q 执行下列程序段后的输出结果是（\quad）。
\begin{lstlisting}[language={[custom]C}]
int i = 3, j;
j = i++ + 2;
printf("%d %d", i, j);
\end{lstlisting}
\choicesline{A. 4 5}{B. 4 4}{C. 3 5}{D. 5 5}

\q 执行下列程序段后，变量 \code{s} 的值是（\quad）。
\begin{lstlisting}[language={[custom]C}]
int i = 0, s = 0;
do {
    s += ++i;
} while (i < 4);
printf("%d", s);
\end{lstlisting}
\choicesline{A. 6}{B. 10}{C. 15}{D. 3}

\q 下列程序段的输出结果是（\quad）。
\begin{lstlisting}[language={[custom]C}]
int a[3][3] = {{1, 2, 3}, {4, 5, 6}, {7, 8, 9}};
int i, s = 0;
for (i = 0; i < 3; i++)
    s += a[i][2 - i];
printf("%d", s);
\end{lstlisting}
\choicesline{A. 15}{B. 9}{C. 12}{D. 18}

\q 下列程序段的输出结果是（\quad）。
\begin{lstlisting}[language={[custom]C}]
char s[] = "abcdef";
char *p = s + 1;
printf("%c%s", *p, p + 2);
\end{lstlisting}
\choicesline{A. bdef}{B. bcdef}{C. bcd}{D. cdef}

\q 下列程序的输出结果是（\quad）。
\begin{lstlisting}[language={[custom]C}]
#include <stdio.h>
int f(int n) {
    if (n <= 1) return 1;
    return f(n - 1) + f(n - 2);
}
int main(void) {
    printf("%d", f(5));
    return 0;
}
\end{lstlisting}
\choicesline{A. 5}{B. 8}{C. 13}{D. 3}

\q 下列程序的输出结果是（\quad）。
\begin{lstlisting}[language={[custom]C}]
#include <stdio.h>
int main(void) {
    int a = 10, b = 20, *p = &a, *q = &b;
    *p = *q;
    q = p;
    *q = 30;
    printf("%d %d", a, b);
    return 0;
}
\end{lstlisting}
\choicesline{A. 30 20}{B. 20 30}{C. 30 30}{D. 20 20}

\q 下列程序的输出结果是（\quad）。
\begin{lstlisting}[language={[custom]C}]
#include <stdio.h>
struct node { int x; struct node *next; };
int main(void) {
    struct node a = {1, NULL}, b = {2, NULL};
    a.next = &b;
    printf("%d", a.next->x + a.x);
    return 0;
}
\end{lstlisting}
\choicesline{A. 1}{B. 2}{C. 3}{D. 编译出错}

\q 关于文件操作，下列说法中正确的是（\quad）。
\choices{A. \code{fopen} 打开文件失败时返回 \code{NULL}，因此必须对返回值进行判断}{B. 调用 \code{fclose} 关闭文件之后，还可以继续用同一个文件指针向该文件写入数据}{C. \code{fgetc} 函数一次可以从文件中读入一个字符串}{D. 以 \code{"w"} 方式打开一个已存在的文件时，原文件的内容会被保留}

\q 下列库函数中，用于把数字字符串转换为整数的是（\quad）。
\choicesline{A. \code{atoi}}{B. \code{itoa}}{C. \code{atof}}{D. \code{sprintf}}
\end{qgroup}

\qtypename{二、填空题}{5 题，每题 2 分，共 10 分}
\begin{qgroup}
\q 若有定义 \code{int x = 10;}，则表达式 \code{x > 0 ? x > 5 ? 1 : 2 : 3} 的值是 \fillblank{}。

\q 下面程序用辗转相除法求两个正整数 \code{m}、\code{n} 的最大公约数，请填空使程序完整正确。
\begin{lstlisting}[language={[custom]C}]
#include <stdio.h>
int main(void) {
    int m, n, r;
    scanf("%d%d", &m, &n);
    while (            ) {
        r = m % n;
        m = n;
        n = r;
    }
    printf("%d\n", m);
    return 0;
}
\end{lstlisting}

\q 若有定义 \code{char s[10] = "abc";}，则 \code{s[3]} 的值是 \fillblank{}（用字符常量表示），\code{sizeof(s)} 的值是 \fillblank{}。

\q 若有定义 \code{int a[5] = \{1, 2, 3, 4, 5\}; int *p = a + 3;}，则 \code{*p} 的值是 \fillblank{}，\code{p - a} 的值是 \fillblank{}。

\q 下面程序把 1 到 5 写入文本文件 \code{num.txt} 再读回并显示，请填空。
\begin{lstlisting}[language={[custom]C}]
#include <stdio.h>
int main(void) {
    FILE *fp;
    int i, x;
    fp = fopen("num.txt", "w");
    for (i = 1; i <= 5; i++) fprintf(fp, "%d ", i);
    (            );            /* 先关闭文件，再以读方式重新打开 */
    fp = fopen("num.txt", "r");
    while (fscanf(fp, "%d", &x) == 1) printf("%d ", x);
    fclose(fp);
    return 0;
}
\end{lstlisting}
\end{qgroup}

\qtypename{三、程序阅读题}{2 题，每题 5 分，共 10 分}
\begin{qgroup}
\q （5 分）写出下列程序的运行结果。
\begin{lstlisting}[language={[custom]C}]
#include <stdio.h>
int f(int n) {
    static int s = 0;
    if (n <= 0) return s;
    s += n;
    return f(n - 1);
}
int main(void) {
    int a = f(4);
    int b = f(2);
    printf("%d %d\n", a, b);
    return 0;
}
\end{lstlisting}
\anslines{1.5cm}

\q （5 分）下列程序用于计算数组元素的累加和与平均值，其中有 3 处错误，请指出错误所在行并改正。
\begin{lstlisting}[language={[custom]C}]
#include <stdio.h>
#define N 5
int main(void)
{
    int a[N] = {1, 2, 3, 4, 6};
    int i, sum = 0;
    float avg;
    for (i = 0; i <= N; i++)
        sum += a[i];
    avg = sum / N;
    printf("sum=%d, avg=%f\n", sum, avg)
    return 0;
}
\end{lstlisting}
\anslines{2.0cm}
\end{qgroup}

\qtypename{四、程序设计题}{2 题，每题 5 分，共 10 分}
\begin{qgroup}
\q （5 分）编写函数 \code{void reverse(char *s)}，把字符串 \code{s} 就地逆序存放（不允许使用额外的数组）。再编写主函数，输入一个字符串，调用该函数后输出逆序后的结果。请写出完整的 C 程序。
\anslines{5.0cm}

\q （5 分）某班有 5 名学生，每名学生的数据包括学号（\code{int}）、姓名（字符串）和 3 门课程的成绩（\code{float}）。要求用结构体数组存放这些数据，从键盘读入后计算每名学生的平均成绩，并输出平均成绩最高的学生的姓名和平均成绩。请写出完整的 C 程序。
\anslines{5.0cm}
\end{qgroup}

\end{document}
